\documentclass[main.tex]{subfiles}
\begin{document}

\begin{table}[H]
    \centering
    \begin{tabular}{p{2cm} p{5.5cm} p{5.5cm}}
         & \textbf{MVC} & \textbf{Observer-Pattern} \\     
       \textbf{Paradigma}  &  Architektur-Muster & Entwurfsmuster \\ 
       \textbf{Zweck} & Strukturiert Software in drei Schichten & Ermöglicht Benachrichtigungen bei Änderungen \\
       \textbf{Kopplung} & Lose Kopplung zwischen Model, View, Controller & Sehr lose Kopplung zwischen Subjekt und Beobachtern \\
       \textbf{Verwendung} & Wird für die gesamte Softwarestruktur verwendet & Wird meist zur Implementierung von Benachrichtigungen genutzt \\
       \textbf{Beziehung zueinander} & MVC kann das Observer-Pattern für die Verbindung zwischen Model und View nutzen & Kann in MVC, aber auch unabhängig verwendet werden  \\ 
       \textbf{Vorteile} & Klare Trennung von Verantwortlichkeiten, Erhöhte Wiederverwendbarkeit von Komponenten, Verbesserte Wartbarkeit und Testbarkeit, Ermöglicht parallele Entwicklung von UI und Logik & Ermöglicht lose Kopplung zwischen Subjekt und Beobachtern, Mehrere Views können automatisch aktualisiert werden, Erhöht die Flexibilität bei Änderungen, Einfach zu implementieren und modular erweiterbar \\
       \textbf{Nachteile} & Erhöhter Implementierungsaufwand, komplex, Erfordert oft zusätzliches Design für Datenbindung und Zustandsmanagement & Kann zu Performance-Problemen führen, wenn zu viele Beobachter registriert sind, Zyklische Abhängigkeiten sind möglich, Schwieriger zu debuggen, wenn viele Events gleichzeitig verarbeitet werden   
    \end{tabular}
    \label{tab:my_label}
\end{table}

\subsection{Fragen}
\subsubsection*{Welche Bestandteile hat MVC?}
\begin{itemize}
    \item Model
    \item View
    \item Controller
\end{itemize}
\subsubsection*{Zeichnen von MVC-Interaktionen}
Aktion im View -> View informiert Controller -> Controller verarbeitet Eingabe und aktualisiert Model -> Model ändert Zustand und benachrichtigt View -> View aktualisiert sich 
\subsubsection*{Vorteil(e) des Observer-Patterns?}
\begin{itemize}
    \item Entkopplung
    \item Flexibilität
    \item Einfache Erweiterbarkeit
\end{itemize}
\subsubsection*{Vorteile von MVC?}
\begin{itemize}
    \item Bessere Wartbarkeit
    \item Wiederverwendbarkeit
    \item Bessere Testbarkeit
    \item Flexibilität
\end{itemize}
\subsubsection*{Warum nutzt man Architekturmuster?}
\begin{itemize}
    \item Strukturierung
    \item Wiederverwendbarkeit
    \item Wartbarkeit
    \item Testbarkeit
    \item Skalierbarkeit
\end{itemize}
\subsubsection*{Wie sind Model und View gekoppelt?}
Gar nicht, das Model sollte keine direkte Abhängigkeit zur View haben. Aber der View ist der Beobachter des Models.  
\subsubsection*{Warum sollte man keine View-Referenzen im Model speichern?}
\begin{itemize}
    \item Verletzung des MVC-Prinzips
    \item Schlechte Wartbarkeit
    \item Schwierige Testbarkeit
    \item Erweiterbarkeit wird erschwert
\end{itemize}
\end{document}